% Chapter 27: The Kakeya Conjecture
\let\LocalMacro\relax

\chapter{Kakeya Conjecture}

\section{The Problem}

\subsection{Informal}
Imagine trying to rotate a needle through every possible direction inside a region of the plane. How small can that region be? In 1917, Soichi Kakeya asked this deceptively simple question.

In 1919, Abram Besicovitch shocked mathematicians by showing the answer is \emph{infinitely small}: you can build a shape that contains a line segment pointing in \emph{every} direction yet occupies essentially zero area. Since area alone can't capture the size of such strange shapes, mathematicians turned to a more refined measure---the \emph{dimension} of the set---leading to a deep conjecture about how ``thick'' these sets must really be:

\subsection{Formal}
\begin{conjecture}[Kakeya Conjecture]
Any Kakeya set in $\mathbb{R}^n$ must have Hausdorff dimension and Minkowski dimension equal to $n$.
\end{conjecture}

Despite its geometric statement, the Kakeya conjecture is deeply connected to harmonic analysis, specifically to wave equations, Fourier multipliers, and the restriction conjecture.

\subsection{Historical Context}
The Kakeya conjecture is one of the most famous unsolved problems in analysis. While it is settled for $n=2$, it remains open in higher dimensions ($n \geq 3$), representing a key barrier to understanding multi-directional wave propagation. Nets Katz of Rice University called the 3D progress the solution of the ``holy grail'' problem in harmonic analysis.

\subsection{What Was Known}
Besicovitch proved in 1919 that Kakeya sets can have zero area, which was already shocking. The question then became: how small can their dimension be? In two dimensions, Davies proved in 1971 that Kakeya sets must have dimension 2. In higher dimensions, the conjecture remained wide open. Finite field analogues were proved by Dvir in 2008, but the real case resisted all attacks.

\subsection{Why It Matters}
The Kakeya conjecture is about the size of the smallest region that can contain a needle pointing in every direction. It connects to harmonic analysis (the study of waves), partial differential equations (the study of fluid flow), and even quantum physics. Understanding how thick Kakeya sets must be has implications for the behavior of solutions to the wave equation.

\subsection{Simplified Version}
The core idea is that containing a line segment in every direction forces a set to be ``thick'' in a precise geometric sense. In two dimensions, this is already proved: any Kakeya set in $\mathbb{R}^2$ must have Hausdorff dimension 2, meaning it is as large as a full-dimensional region even though it can have zero area. The conjecture asserts that this holds in all dimensions---Kakeya sets in $\mathbb{R}^n$ must have dimension $n$. The finite field analogue, proved by Dvir in 2008 using the polynomial method, was a crucial stepping stone toward the real case.

\section{Mathematical Theory}

\subsection{Prerequisites}
This chapter assumes familiarity with:
\begin{itemize}
\item Basic real analysis (Lebesgue measure, open/closed sets, limits)
\item Geometric measure theory (Hausdorff dimension)
\item Finite fields and polynomials (roots, finite vector spaces)
\item Harmonic analysis (Fourier transform, restriction operators) --- helpful but not essential
\end{itemize}

\subsection{Besicovitch's Construction}
Besicovitch's proof uses a geometric decomposition called ``panning.'' By dividing a triangle into smaller triangles, shifting them, and overlapping them, one can concentrate the area into arbitrarily small regions while preserving segments in all directions.

\subsection{Finite Field Analogue}
In 1999, Thomas Wolff proposed studying the Kakeya conjecture in the simplified setting of vector spaces over finite fields $\mathbb{F}_q^n$ to bypass real-variable complications \cite{wolff1999}.

\begin{definition}[Finite Field Kakeya Set]
A subset $K \subset \mathbb{F}_q^n$ is a Kakeya set if for every direction $y \in \mathbb{F}_q^n \setminus \{0\}$, there exists $x \in \mathbb{F}_q^n$ such that $\{x + t y : t \in \mathbb{F}_q\} \subset K$.
\end{definition}

Wolff conjectured that any Kakeya set in $\mathbb{F}_q^n$ satisfies $|K| \geq C_n q^n$.

\section{The Solution}

\subsection{What Was Missing}
The Kakeya conjecture---that Kakeya sets in $\mathbb{R}^n$ must have Hausdorff dimension $n$---was settled for $n=2$ (Davies, 1971) but remained wide open in $\mathbb{R}^3$ and higher. The best known lower bound in $\mathbb{R}^3$ was stuck below 2.5 for many years. What was missing was a way to break past this dimension barrier: existing Fourier-analytic and combinatorial methods could not extract dimensional gains beyond roughly $n/2 + \epsilon$, and the polynomial method that had been spectacularly successful for the finite field analogue (Dvir, 2008) had no direct translation to the real setting.

\subsection{Why Existing Methods Failed}
Fourier-analytic methods---restriction theory, decoupling, and endpoint estimates---gave incremental improvements but always hit the same ceiling around dimension $2.5$ in $\mathbb{R}^3$. The polynomial method that Dvir used over $\mathbb{F}_q$ relies on the fact that a polynomial of degree $<q$ vanishing on all points of a line must vanish identically; this argument fails over $\mathbb{R}$ because lines have infinitely many points and there is no degree bound. Conversely, incidence geometry methods (Guth--Katz) gave powerful bounds for discrete problems but required significant adaptation to handle the continuous, measure-theoretic setting of Kakeya sets.

\subsection{What Was Novel}
Wang and Zahl's 2025 breakthrough in $\mathbb{R}^3$ combined three innovations adapted from the discrete to the continuous: (1)~\textbf{polynomial partitioning} (Guth--Katz, 2010) to decompose space into cells each containing relatively few Kakeya lines; (2)~\textbf{linear programming bounds} that optimize over families of line configurations to extract dimensional lower bounds; and (3)~\textbf{new incidence geometry} giving sharp bounds on point-line incidences in three dimensions. Together these tools broke the 2.5 barrier, proving that Kakeya sets in $\mathbb{R}^3$ have Hausdorff dimension at least $2.5 + \epsilon$ for an absolute $\epsilon > 0$.

\subsection{Student-Friendly Explanation}
A Kakeya set is a region of space that contains a line segment in every direction---like a room that has a needle pointing every possible way. You might expect such a room to be large, but Besicovitch showed it can have zero volume. The question is: can it also be ``thin'' in a more refined sense (dimension less than $n$)? In two dimensions the answer is no (every Kakeya set has dimension 2), but in three dimensions this was stuck for decades. Wang and Zahl's breakthrough used algebraic cutting: they sliced space with the zero set of a polynomial, counted how many needle-lines fall in each piece, and used an optimization argument to show the needles must pack together thickly enough to force dimension above 2.5.

The Kakeya conjecture has seen two major breakthroughs in recent years.

\subsection{Dvir's Finite Field Proof (2008)}
In 2008, Zeev Dvir proved Wolff's finite field conjecture in a paper that was only four pages long, using the \emph{polynomial method} \cite{dvir2009}:

\begin{theorem}[Dvir, 2008 \cite{dvir2009}]
Any Kakeya set $K \subset \mathbb{F}_q^n$ satisfies:
\begin{equation}
|K| \geq \frac{1}{n!} q^n
\end{equation}
\end{theorem}

The proof is elegant: if $|K| < \binom{n+q-1}{n}$, there exists a non-zero polynomial of degree $<q$ vanishing on $K$. By complete induction along line directions, the polynomial's highest degree homogeneous part must vanish on all of $\mathbb{F}_q^n$, which is impossible.

\subsection{The 3D Real Breakthrough: Wang--Zahl (2023--2025)}

The polynomial method had been spectacularly successful for the finite field Kakeya problem, but extending it to the real setting $\mathbb{R}^n$ had resisted all attempts for decades. The breakthrough came from Hong Wang (NYU) and Joshua Zahl (UBC).

\begin{historicalbox}
Hong Wang (born 1992, China) brought deep expertise in polynomial partitioning and geometric measure theory. Joshua Zahl (UBC) contributed experience in incidence geometry and harmonic analysis. Their collaboration spanned several years and required new ideas at each stage.
\end{historicalbox}

\begin{theorem}[Wang--Zahl, 2025 \cite{wangzahl2025}]
Any Kakeya set $K \subset \mathbb{R}^3$ has Hausdorff dimension at least $2.5 + \epsilon$ for an absolute constant $\epsilon > 0$.
\end{theorem}

Their proof combined three major innovations:

\begin{enumerate}
\item \textbf{Polynomial partitioning}: A technique from algebraic geometry (introduced by Guth and Katz in 2010) that decomposes space using the zero sets of polynomials. Given a Kakeya set $K \subset \mathbb{R}^3$, construct a polynomial $P$ that vanishes on a carefully chosen subset of the lines in $K$. The zero set of $P$ decomposes space into cells, each containing relatively few lines.
\item \textbf{Linear programming bounds}: Optimizing over families of line configurations to extract dimensional lower bounds. This technique converts geometric statements about Kakeya sets into optimization problems that can be solved systematically.
\item \textbf{New incidence geometry}: Novel bounds on the number of incidences between points and lines in three dimensions, extending the work of Guth and Katz.
\end{enumerate}

\begin{highlightbox}
The Wang--Zahl result broke the long-standing ``2.5 dimension barrier'' in $\mathbb{R}^3$. The full Kakeya conjecture---that Kakeya sets in $\mathbb{R}^n$ have Hausdorff dimension $n$---remains open for $n \geq 3$, but the gap between the best lower bound and the conjectured value has narrowed dramatically.
\end{highlightbox}

\subsection{Hong Wang and the Fields Medal}

In 2026, Hong Wang was awarded the Fields Medal at the International Congress of Mathematicians in Philadelphia, becoming the third woman to receive the prize (after Maryam Mirzakhani in 2014 and Maryna Viazovska in 2022). Her work demonstrated that algebraic techniques are powerful enough to attack real-variable problems that had resisted Fourier-analytic methods for decades.

\section{Impact}
\begin{itemize}
\item \textbf{Harmonic Analysis}: Directly bounds the restriction operator for the Fourier transform, influencing local smoothing estimates for wave equations. The Wang--Zahl bound limits how concentrated energy can become during wave propagation.
\item \textbf{Additive Combinatorics}: The finite field polynomial method has been adapted to prove the joints conjecture and other incidence geometry problems. The incidence geometry techniques have applications to sum-product estimates.
\item \textbf{Partial Differential Equations}: Provides bounds on the size of singularity sets for dispersive wave propagation.
\item \textbf{Geometric Measure Theory}: The polynomial method developed for Kakeya has been applied to Falconer's distance set problem and other classical questions.
\end{itemize}

\section{Exercises}

\begin{exercise}[Basic]
Show that a Kakeya set in $\mathbb{R}^2$ containing segments in all directions can be constructed with measure 0.
\end{exercise}

\begin{exercise}[Intermediate]
Write down the details of Zeev Dvir's finite field Kakeya set proof. Show why the degree of the vanishing polynomial must be less than $q$.
\end{exercise}

\begin{exercise}[Challenge]
Explain why Dvir's algebraic proof fails in the setting of real numbers $\mathbb{R}^n$. What is the main difference between lines in $\mathbb{F}_q^n$ and $\mathbb{R}^n$?
\end{exercise}

\begin{exercise}[Computational]
Write a Python/SageMath script to explore a concept from this chapter. See the appendix for starter code.
\end{exercise}

\section{Timeline}

\begin{tabularx}{\linewidth}{lX}
\toprule
\textbf{Year} & \textbf{Event} \\ \midrule
1917 & Kakeya poses the needle turning problem \cite{kakeya1917} \\
1919 & Besicovitch constructs a Kakeya set of measure zero \cite{besicovitch1919} \\
1971 & C\'ordoba proves the Kakeya maximal function conjecture for $n=2$ \\
1999 & Wolff formulates the finite field Kakeya conjecture \cite{wolff1999} \\
2003 & Wolff improves the lower bound on Kakeya dimension in $\mathbb{R}^3$ \\
2008 & Dvir proves the finite field Kakeya conjecture using the polynomial method \cite{dvir2009} \\
2023--2025 & Wang and Zahl break the 2.5-dimension barrier in $\mathbb{R}^3$ \cite{wangzahl2024} \\
2026 & Hong Wang awarded the Fields Medal at ICM Philadelphia \\
\bottomrule
\end{tabularx}

\section{Citations}

\begin{enumerate}
\item Kakeya, S. (1917). ``Some problems on maximum and minimum regarding ovals.'' Tohoku Mathematical Journal.
\item Besicovitch, A. S. (1919). ``Sur deux questions de geometrie.'' Bulletin de l'Academie des Sciences de Russie.
\item Wolff, T. (1999). ``Recent work on restriction and Kakeya conjectures.'' IAS Lectures.
\item Dvir, Z. (2009). ``On the size of Kakeya sets in finite fields.'' Journal of the AMS, 23(4), 1107--1119.
\item Guth, L., \& Katz, N. (2015). ``On the Erd\H{o}s distinct distances problem in the plane.'' Annals of Mathematics, 181(1), 155--190.
\item Wang, H., \& Zahl, J. (2025). ``The Kakeya conjecture in three dimensions.'' arXiv:2502.17655.
\item International Mathematical Union. (2026). ``Fields Medal citation: Hong Wang.'' ICM 2026, Philadelphia.
\end{enumerate}
